% =====================================================================
%  Anatomía de un LLM — Fase II — Módulo 9
%  Redes neuronales
%  Compilar: latexmk -r ../latexmkrc M09-redes-neuronales.tex
% =====================================================================
\documentclass[11pt,a4paper,twoside,openany]{book}
\usepackage{../estilo-apuntes}

\newcommand{\ApuntesCurso}{Anatomía de un LLM}
\newcommand{\ApuntesModulo}{Módulo 9}

% Notación propia del módulo
\newcommand{\ReLU}{\operatorname{ReLU}}
\newcommand{\GELU}{\operatorname{GELU}}
\newcommand{\SiLU}{\operatorname{SiLU}}
\newcommand{\Var}{\operatorname{Var}}
\newcommand{\ind}{\mathbb{1}}
\newcommand{\FFN}{\operatorname{FFN}}
\newcommand{\nin}{n_{\mathrm{in}}}
\newcommand{\nout}{n_{\mathrm{out}}}
\makeatletter\g@addto@macro\UrlBreaks{\do\-}\makeatother

\hypersetup{pdftitle={Módulo 9. Redes neuronales (resumen)},pdfauthor={Anatomía de un LLM}}

\makeatletter
\newlabel{ej:and}{{1.3}{1}{La puerta AND}{ejemplo.3}{}}
\newlabel{ej:and@cref}{{[ejemplo][3][]1.3}{[1][1][]1}}
\newlabel{ej:rosenblatt}{{1.6}{2}{Aprendizaje de AND}{ejemplo.6}{}}
\newlabel{ej:rosenblatt@cref}{{[ejemplo][6][]1.6}{[1][2][]2}}
\newlabel{fig:xor}{{1.1}{3}{Red con una capa oculta de dos unidades ReLU que calcula XOR}{figure.1.1}{}}
\newlabel{fig:xor@cref}{{[figure][1][1]1.1}{[1][3][]3}}
\newlabel{fig:mlp}{{1.2}{3}{Perceptrón multicapa con una capa oculta de anchura $H$ y salida softmax}{figure.1.2}{}}
\newlabel{fig:mlp@cref}{{[figure][2][1]1.2}{[1][3][]3}}
\newlabel{ej:mnist}{{1.12}{4}{}{ejemplo.12}{}}
\newlabel{ej:mnist@cref}{{[ejemplo][12][]1.12}{[1][4][]4}}
\newlabel{eje:1-lineal}{{1.1}{4}{}{ejercicio.1.1}{}}
\newlabel{eje:1-lineal@cref}{{[ejercicio][1][1]1.1}{[1][4][]4}}
\newlabel{eje:1-xor}{{1.2}{4}{}{ejercicio.1.2}{}}
\newlabel{eje:1-xor@cref}{{[ejercicio][2][1]1.2}{[1][4][]4}}
\newlabel{eje:1-param}{{1.3}{5}{}{ejercicio.1.3}{}}
\newlabel{eje:1-param@cref}{{[ejercicio][3][1]1.3}{[1][5][]5}}
\newlabel{eje:1-cs231n}{{1.4}{5}{}{ejercicio.1.4}{}}
\newlabel{eje:1-cs231n@cref}{{[ejercicio][4][1]1.4}{[1][5][]5}}
\newlabel{eje:1-or}{{1.5}{5}{}{ejercicio.1.5}{}}
\newlabel{eje:1-or@cref}{{[ejercicio][5][1]1.5}{[1][5][]5}}
\newlabel{eje:1-novikoff}{{1.6}{5}{}{ejercicio.1.6}{}}
\newlabel{eje:1-novikoff@cref}{{[ejercicio][6][1]1.6}{[1][5][]5}}
\newlabel{fig:activaciones}{{2.1}{7}{ReLU, GELU y SiLU en $[-3,3]$. Las dos activaciones suaves dejan pasar una pequeña señal negativa}{figure.2.1}{}}
\newlabel{fig:activaciones@cref}{{[figure][1][2]2.1}{[1][7][]7}}
\newlabel{ej:swiglu}{{2.8}{8}{}{ejemplo.8}{}}
\newlabel{ej:swiglu@cref}{{[ejemplo][8][]2.8}{[1][8][]8}}
\newlabel{eje:2-sigmoide}{{2.1}{9}{}{ejercicio.2.1}{}}
\newlabel{eje:2-sigmoide@cref}{{[ejercicio][1][2]2.1}{[1][9][]9}}
\newlabel{eje:2-anchura}{{2.2}{9}{}{ejercicio.2.2}{}}
\newlabel{eje:2-anchura@cref}{{[ejercicio][2][2]2.2}{[1][9][]9}}
\newlabel{eje:2-tanh}{{2.3}{9}{}{ejercicio.2.3}{}}
\newlabel{eje:2-tanh@cref}{{[ejercicio][3][2]2.3}{[1][9][]9}}
\newlabel{eje:2-swish}{{2.4}{9}{}{ejercicio.2.4}{}}
\newlabel{eje:2-swish@cref}{{[ejercicio][4][2]2.4}{[1][9][]9}}
\newlabel{eje:2-desvanece}{{2.5}{9}{}{ejercicio.2.5}{}}
\newlabel{eje:2-desvanece@cref}{{[ejercicio][5][2]2.5}{[1][9][]9}}
\newlabel{eje:2-swiglu}{{2.6}{9}{}{ejercicio.2.6}{}}
\newlabel{eje:2-swiglu@cref}{{[ejercicio][6][2]2.6}{[1][9][]9}}
\newlabel{eje:2-swish-relu}{{2.7}{9}{}{ejercicio.2.7}{}}
\newlabel{eje:2-swish-relu@cref}{{[ejercicio][7][2]2.7}{[1][9][]9}}
\newlabel{ej:cuadrado}{{3.5}{11}{Aproximación de $t^{2}$}{ejemplo.5}{}}
\newlabel{ej:cuadrado@cref}{{[ejemplo][5][]3.5}{[1][11][]11}}
\newlabel{eje:3-existencia}{{3.1}{13}{}{ejercicio.3.1}{}}
\newlabel{eje:3-existencia@cref}{{[ejercicio][1][3]3.1}{[1][12][]13}}
\newlabel{eje:3-codos}{{3.2}{13}{}{ejercicio.3.2}{}}
\newlabel{eje:3-codos@cref}{{[ejercicio][2][3]3.2}{[1][13][]13}}
\newlabel{eje:3-cuadrado}{{3.3}{13}{}{ejercicio.3.3}{}}
\newlabel{eje:3-cuadrado@cref}{{[ejercicio][3][3]3.3}{[1][13][]13}}
\newlabel{eje:3-sierra}{{3.4}{13}{}{ejercicio.3.4}{}}
\newlabel{eje:3-sierra@cref}{{[ejercicio][4][3]3.4}{[1][13][]13}}
\newlabel{eje:3-abs-max}{{3.5}{13}{}{ejercicio.3.5}{}}
\newlabel{eje:3-abs-max@cref}{{[ejercicio][5][3]3.5}{[1][13][]13}}
\newlabel{eje:3-cpwl}{{3.6}{13}{}{ejercicio.3.6}{}}
\newlabel{eje:3-cpwl@cref}{{[ejercicio][6][3]3.6}{[1][13][]13}}
\newlabel{ej:xavier}{{4.7}{15}{}{ejemplo.7}{}}
\newlabel{ej:xavier@cref}{{[ejemplo][7][]4.7}{[1][15][]15}}
\newlabel{ej:he}{{4.10}{16}{}{ejemplo.10}{}}
\newlabel{ej:he@cref}{{[ejemplo][10][]4.10}{[1][16][]16}}
\newlabel{eje:4-simetria}{{4.1}{17}{}{ejercicio.4.1}{}}
\newlabel{eje:4-simetria@cref}{{[ejercicio][1][4]4.1}{[1][17][]17}}
\newlabel{eje:4-xavier}{{4.2}{17}{}{ejercicio.4.2}{}}
\newlabel{eje:4-xavier@cref}{{[ejercicio][2][4]4.2}{[1][17][]17}}
\newlabel{eje:4-uniforme}{{4.3}{17}{}{ejercicio.4.3}{}}
\newlabel{eje:4-uniforme@cref}{{[ejercicio][3][4]4.3}{[1][17][]17}}
\newlabel{eje:4-decaimiento}{{4.4}{17}{}{ejercicio.4.4}{}}
\newlabel{eje:4-decaimiento@cref}{{[ejercicio][4][4]4.4}{[1][17][]17}}
\newlabel{eje:4-he}{{4.5}{17}{}{ejercicio.4.5}{}}
\newlabel{eje:4-he@cref}{{[ejercicio][5][4]4.5}{[1][17][]17}}
\newlabel{eje:4-residual}{{4.6}{17}{}{ejercicio.4.6}{}}
\newlabel{eje:4-residual@cref}{{[ejercicio][6][4]4.6}{[1][17][]17}}
\newlabel{fig:grafo}{{5.1}{20}{Grafo computacional del MLP del \cref {teo:backprop}. Las flechas continuas son la pasada hacia delante; las discontinuas, la pasada hacia atrás}{figure.5.1}{}}
\newlabel{fig:grafo@cref}{{[figure][1][5]5.1}{[1][20][]20}}
\newlabel{ej:backprop}{{5.10}{20}{Red $2\to 3\to 3$ con ReLU}{ejemplo.10}{}}
\newlabel{ej:backprop@cref}{{[ejemplo][10][]5.10}{[1][20][]20}}
\newlabel{eje:5-softmax}{{5.1}{21}{}{ejercicio.5.1}{}}
\newlabel{eje:5-softmax@cref}{{[ejercicio][1][5]5.1}{[1][21][]21}}
\newlabel{eje:5-delta}{{5.2}{21}{}{ejercicio.5.2}{}}
\newlabel{eje:5-delta@cref}{{[ejercicio][2][5]5.2}{[1][21][]21}}
\newlabel{eje:5-y0}{{5.3}{22}{}{ejercicio.5.3}{}}
\newlabel{eje:5-y0@cref}{{[ejercicio][3][5]5.3}{[1][21][]22}}
\newlabel{eje:5-cs224n}{{5.4}{22}{}{ejercicio.5.4}{}}
\newlabel{eje:5-cs224n@cref}{{[ejercicio][4][5]5.4}{[1][22][]22}}
\newlabel{eje:5-grafo}{{5.5}{22}{}{ejercicio.5.5}{}}
\newlabel{eje:5-grafo@cref}{{[ejercicio][5][5]5.5}{[1][22][]22}}
\newlabel{eje:5-suma}{{5.6}{22}{}{ejercicio.5.6}{}}
\newlabel{eje:5-suma@cref}{{[ejercicio][6][5]5.6}{[1][22][]22}}
\newlabel{eje:5-swiglu}{{5.7}{22}{}{ejercicio.5.7}{}}
\newlabel{eje:5-swiglu@cref}{{[ejercicio][7][5]5.7}{[1][22][]22}}
\newlabel{fig:bengio}{{6.1}{24}{Modelo de lenguaje neuronal de la práctica, con contexto de tres tokens}{figure.6.1}{}}
\newlabel{fig:bengio@cref}{{[figure][1][6]6.1}{[1][24][]24}}
\newlabel{ej:practica}{{6.5}{24}{El modelo de la práctica}{ejemplo.5}{}}
\newlabel{ej:practica@cref}{{[ejemplo][5][]6.5}{[1][24][]24}}
\newlabel{eje:6-sobreajuste}{{6.1}{26}{}{ejercicio.6.1}{}}
\newlabel{eje:6-sobreajuste@cref}{{[ejercicio][1][6]6.1}{[1][26][]26}}
\newlabel{eje:6-tabla}{{6.2}{26}{}{ejercicio.6.2}{}}
\newlabel{eje:6-tabla@cref}{{[ejercicio][2][6]6.2}{[1][26][]26}}
\newlabel{eje:6-param}{{6.3}{26}{}{ejercicio.6.3}{}}
\newlabel{eje:6-param@cref}{{[ejercicio][3][6]6.3}{[1][26][]26}}
\newlabel{eje:6-embgrad}{{6.4}{26}{}{ejercicio.6.4}{}}
\newlabel{eje:6-embgrad@cref}{{[ejercicio][4][6]6.4}{[1][26][]26}}
\newlabel{eje:6-bigrama}{{6.5}{26}{}{ejercicio.6.5}{}}
\newlabel{eje:6-bigrama@cref}{{[ejercicio][5][6]6.5}{[1][26][]26}}
\newlabel{ap:soluciones}{{A}{27}{Soluciones de los ejercicios}{appendix.A}{}}
\newlabel{ap:soluciones@cref}{{[appendix][1][2147483647]A}{[1][27][]27}}
\newlabel{sol:eje:1-lineal}{{A.1}{27}{Tema 1}{section*.3}{}}
\newlabel{sol:eje:1-lineal@cref}{{[subappendix][1][2147483647,1]A.1}{[1][27][]27}}
\newlabel{sol:eje:1-xor}{{A.1}{27}{Tema 1}{section*.4}{}}
\newlabel{sol:eje:1-xor@cref}{{[subappendix][1][2147483647,1]A.1}{[1][27][]27}}
\newlabel{sol:eje:1-param}{{A.1}{28}{Tema 1}{section*.5}{}}
\newlabel{sol:eje:1-param@cref}{{[subappendix][1][2147483647,1]A.1}{[1][28][]28}}
\newlabel{sol:eje:1-cs231n}{{A.1}{28}{Tema 1}{section*.6}{}}
\newlabel{sol:eje:1-cs231n@cref}{{[subappendix][1][2147483647,1]A.1}{[1][28][]28}}
\newlabel{sol:eje:1-or}{{A.1}{28}{Tema 1}{section*.7}{}}
\newlabel{sol:eje:1-or@cref}{{[subappendix][1][2147483647,1]A.1}{[1][28][]28}}
\newlabel{sol:eje:1-novikoff}{{A.1}{29}{Tema 1}{section*.8}{}}
\newlabel{sol:eje:1-novikoff@cref}{{[subappendix][1][2147483647,1]A.1}{[1][29][]29}}
\newlabel{sol:eje:2-sigmoide}{{A.2}{29}{Tema 2}{section*.9}{}}
\newlabel{sol:eje:2-sigmoide@cref}{{[subappendix][2][2147483647,1]A.2}{[1][29][]29}}
\newlabel{sol:eje:2-anchura}{{A.2}{30}{Tema 2}{section*.10}{}}
\newlabel{sol:eje:2-anchura@cref}{{[subappendix][2][2147483647,1]A.2}{[1][30][]30}}
\newlabel{sol:eje:2-tanh}{{A.2}{30}{Tema 2}{section*.11}{}}
\newlabel{sol:eje:2-tanh@cref}{{[subappendix][2][2147483647,1]A.2}{[1][30][]30}}
\newlabel{sol:eje:2-swish}{{A.2}{30}{Tema 2}{section*.12}{}}
\newlabel{sol:eje:2-swish@cref}{{[subappendix][2][2147483647,1]A.2}{[1][30][]30}}
\newlabel{sol:eje:2-desvanece}{{A.2}{31}{Tema 2}{section*.13}{}}
\newlabel{sol:eje:2-desvanece@cref}{{[subappendix][2][2147483647,1]A.2}{[1][31][]31}}
\newlabel{sol:eje:2-swiglu}{{A.2}{31}{Tema 2}{section*.14}{}}
\newlabel{sol:eje:2-swiglu@cref}{{[subappendix][2][2147483647,1]A.2}{[1][31][]31}}
\newlabel{sol:eje:2-swish-relu}{{A.2}{31}{Tema 2}{section*.15}{}}
\newlabel{sol:eje:2-swish-relu@cref}{{[subappendix][2][2147483647,1]A.2}{[1][31][]31}}
\newlabel{sol:eje:3-existencia}{{A.3}{32}{Tema 3}{section*.16}{}}
\newlabel{sol:eje:3-existencia@cref}{{[subappendix][3][2147483647,1]A.3}{[1][32][]32}}
\newlabel{sol:eje:3-codos}{{A.3}{32}{Tema 3}{section*.17}{}}
\newlabel{sol:eje:3-codos@cref}{{[subappendix][3][2147483647,1]A.3}{[1][32][]32}}
\newlabel{sol:eje:3-cuadrado}{{A.3}{33}{Tema 3}{section*.18}{}}
\newlabel{sol:eje:3-cuadrado@cref}{{[subappendix][3][2147483647,1]A.3}{[1][33][]33}}
\newlabel{sol:eje:3-sierra}{{A.3}{33}{Tema 3}{section*.19}{}}
\newlabel{sol:eje:3-sierra@cref}{{[subappendix][3][2147483647,1]A.3}{[1][33][]33}}
\newlabel{sol:eje:3-abs-max}{{A.3}{34}{Tema 3}{section*.20}{}}
\newlabel{sol:eje:3-abs-max@cref}{{[subappendix][3][2147483647,1]A.3}{[1][33][]34}}
\newlabel{sol:eje:3-cpwl}{{A.3}{34}{Tema 3}{section*.21}{}}
\newlabel{sol:eje:3-cpwl@cref}{{[subappendix][3][2147483647,1]A.3}{[1][34][]34}}
\newlabel{sol:eje:4-simetria}{{A.4}{35}{Tema 4}{section*.22}{}}
\newlabel{sol:eje:4-simetria@cref}{{[subappendix][4][2147483647,1]A.4}{[1][34][]35}}
\newlabel{sol:eje:4-xavier}{{A.4}{35}{Tema 4}{section*.23}{}}
\newlabel{sol:eje:4-xavier@cref}{{[subappendix][4][2147483647,1]A.4}{[1][35][]35}}
\newlabel{sol:eje:4-uniforme}{{A.4}{35}{Tema 4}{section*.24}{}}
\newlabel{sol:eje:4-uniforme@cref}{{[subappendix][4][2147483647,1]A.4}{[1][35][]35}}
\newlabel{sol:eje:4-decaimiento}{{A.4}{36}{Tema 4}{section*.25}{}}
\newlabel{sol:eje:4-decaimiento@cref}{{[subappendix][4][2147483647,1]A.4}{[1][36][]36}}
\newlabel{sol:eje:4-he}{{A.4}{36}{Tema 4}{section*.26}{}}
\newlabel{sol:eje:4-he@cref}{{[subappendix][4][2147483647,1]A.4}{[1][36][]36}}
\newlabel{sol:eje:4-residual}{{A.4}{36}{Tema 4}{section*.27}{}}
\newlabel{sol:eje:4-residual@cref}{{[subappendix][4][2147483647,1]A.4}{[1][36][]36}}
\newlabel{sol:eje:5-softmax}{{A.5}{37}{Tema 5}{section*.28}{}}
\newlabel{sol:eje:5-softmax@cref}{{[subappendix][5][2147483647,1]A.5}{[1][37][]37}}
\newlabel{sol:eje:5-delta}{{A.5}{37}{Tema 5}{section*.29}{}}
\newlabel{sol:eje:5-delta@cref}{{[subappendix][5][2147483647,1]A.5}{[1][37][]37}}
\newlabel{sol:eje:5-y0}{{A.5}{38}{Tema 5}{section*.30}{}}
\newlabel{sol:eje:5-y0@cref}{{[subappendix][5][2147483647,1]A.5}{[1][38][]38}}
\newlabel{sol:eje:5-cs224n}{{A.5}{38}{Tema 5}{section*.31}{}}
\newlabel{sol:eje:5-cs224n@cref}{{[subappendix][5][2147483647,1]A.5}{[1][38][]38}}
\newlabel{sol:eje:5-grafo}{{A.5}{39}{Tema 5}{section*.32}{}}
\newlabel{sol:eje:5-grafo@cref}{{[subappendix][5][2147483647,1]A.5}{[1][38][]39}}
\newlabel{sol:eje:5-suma}{{A.5}{39}{Tema 5}{section*.33}{}}
\newlabel{sol:eje:5-suma@cref}{{[subappendix][5][2147483647,1]A.5}{[1][39][]39}}
\newlabel{sol:eje:5-swiglu}{{A.5}{39}{Tema 5}{section*.34}{}}
\newlabel{sol:eje:5-swiglu@cref}{{[subappendix][5][2147483647,1]A.5}{[1][39][]39}}
\newlabel{sol:eje:6-sobreajuste}{{A.6}{40}{Tema 6}{section*.35}{}}
\newlabel{sol:eje:6-sobreajuste@cref}{{[subappendix][6][2147483647,1]A.6}{[1][40][]40}}
\newlabel{sol:eje:6-tabla}{{A.6}{40}{Tema 6}{section*.36}{}}
\newlabel{sol:eje:6-tabla@cref}{{[subappendix][6][2147483647,1]A.6}{[1][40][]40}}
\newlabel{sol:eje:6-param}{{A.6}{41}{Tema 6}{section*.37}{}}
\newlabel{sol:eje:6-param@cref}{{[subappendix][6][2147483647,1]A.6}{[1][40][]41}}
\newlabel{sol:eje:6-embgrad}{{A.6}{41}{Tema 6}{section*.38}{}}
\newlabel{sol:eje:6-embgrad@cref}{{[subappendix][6][2147483647,1]A.6}{[1][41][]41}}
\newlabel{sol:eje:6-bigrama}{{A.6}{41}{Tema 6}{section*.39}{}}
\newlabel{sol:eje:6-bigrama@cref}{{[subappendix][6][2147483647,1]A.6}{[1][41][]41}}
\makeatother

\begin{document}
\frontmatter
\begin{titlepage}
\centering
\vspace*{2.5cm}
{\large\scshape Anatomía de un LLM}\par
\vspace{0.4cm}
{\normalsize Fundamentos matemáticos de los modelos de lenguaje}\par
\vspace{2.2cm}
\rule{\textwidth}{0.8pt}\par\vspace{0.5cm}
{\small\scshape Fase II \textperiodcentered\ Aprendizaje automático y redes neuronales}\par\vspace{0.3cm}
{\Huge\bfseries Módulo 9\par}\vspace{0.3cm}
{\LARGE Redes neuronales\par}
\vspace{0.2cm}
{\large Del perceptrón al modelo de lenguaje neuronal\par}
\vspace{0.5cm}\rule{\textwidth}{0.4pt}\par
\vspace{1.6cm}
{\large Resumen del módulo}\par
\vfill
{\small Edición de octubre de 2026}\par
\end{titlepage}

\chapter*{Sobre este resumen}
\addcontentsline{toc}{chapter}{Sobre este resumen}
\markboth{Sobre este resumen}{}

Este documento reúne lo indispensable de los apuntes del Módulo 9: el texto expositivo, todas las definiciones y todos los resultados (teoremas, proposiciones, lemas y corolarios), las tablas y las observaciones de contenido. Se han suprimido las demostraciones, los ejemplos, las lecturas, las figuras, los ejercicios y sus soluciones. La numeración de temas, secciones, enunciados y ecuaciones es la misma que en los apuntes, de modo que cualquier referencia puede buscarse allí; las que remiten a elementos suprimidos (un ejemplo, una figura, una ecuación de una demostración) conservan su número en los apuntes.

\section*{Notación}
\begin{center}\small
\begin{tabularx}{\textwidth}{@{}l>{\raggedright\arraybackslash}X@{}}
\toprule
Símbolo & Significado \\
\midrule
$x\in\R^{1\times n}$ & vector de entrada, escrito como fila \\
$W\in\R^{n\times m}$, $b\in\R^{1\times m}$ & matriz de pesos y vector de sesgos de una capa $n\to m$ \\
$z=xW+b$, $h=\phi(z)$ & preactivación y activación de una capa \\
$\phi$, $\phi'$ & función de activación, aplicada componente a componente, y su derivada \\
$\sigma$, $\Phi$ & sigmoide logística y función de distribución de la normal estándar \\
$\odot$ & producto componente a componente (de Hadamard) \\
$p=\softmax(z)$, $e_y$ & distribución de salida y vector \emph{one-hot} de la clase $y$ \\
$L=-\ln p_y$ & pérdida de entropía cruzada de un ejemplo \\
$\nabla_{\!W}L$, $\delta$ & gradiente respecto a $W$ (misma forma que $W$) y gradiente respecto a una preactivación \\
$\nin$, $\nout$ & número de entradas y de salidas de una capa \\
$\V$, $K=\abs{\V}$, $C\in\R^{K\times d}$ & vocabulario, su tamaño y matriz de embeddings \\
\bottomrule
\end{tabularx}
\end{center}

A diferencia de los módulos~2 y~3, en este módulo los vectores se escriben como \emph{filas} y las capas actúan por la derecha, $x\mapsto xW+b$. Es la convención de NumPy y PyTorch, la de la práctica del módulo y la de los artículos de Shazeer~\cite{shazeer} y Touvron \emph{et al.}~\cite{touvron}; con un lote de $B$ ejemplos, $x$ pasa a ser una matriz $B\times n$ y las fórmulas no cambian. Las clases y los tokens se numeran desde~0.


\tableofcontents

\mainmatter
% =====================================================================
\chapter{Del perceptrón al perceptrón multicapa}
\label{tema:1}
% =====================================================================

Este tema presenta la unidad básica de una red neuronal, el perceptrón, su regla de aprendizaje y el teorema que garantiza su convergencia, y muestra por qué una sola unidad no basta: la función XOR no es linealmente separable. La solución, añadir una capa oculta, conduce al perceptrón multicapa. La exposición sigue la §6.1 de~\cite{gbc} y las §6.1--6.3 de~\cite{jm3}; la demostración del teorema de convergencia sigue la lección~3 de~\cite{cs4780}, basada en Novikoff~\cite{novikoff}, y el recuento de parámetros, las notas «Neural Networks Part~1» de~\cite{cs231n}.

\section{La unidad y el perceptrón}

\begin{definicion}[Unidad]\label{def:unidad}
Una \emph{unidad} (o \emph{neurona artificial}) con $n$ entradas está determinada por un vector de \emph{pesos} $w\in\R^{n}$, un \emph{sesgo} $b\in\R$ y una función $\phi\colon\R\to\R$, llamada \emph{función de activación}. Para una entrada $x\in\R^{n}$ calcula la \emph{preactivación} $z=w\cdot x+b$ y la \emph{activación} $a=\phi(z)$.
\end{definicion}

\begin{definicion}[Perceptrón]\label{def:perceptron}
El \emph{perceptrón} de Rosenblatt~\cite{rosenblatt} es la unidad con activación escalón:
\[
\hat y=\ind[\,w\cdot x+b>0\,]\in\{0,1\}.
\]
Su \emph{frontera de decisión} es el hiperplano $\{x\in\R^{n}:w\cdot x+b=0\}$.
\end{definicion}

El perceptrón es, por tanto, un clasificador lineal: asigna la clase~1 al semiespacio abierto $\{w\cdot x+b>0\}$ y la clase~0 a su complementario.

\addtocounter{definicion}{1}%

\begin{definicion}[Separabilidad lineal]
Un conjunto finito de pares $(x_i,y_i)$ con $x_i\in\R^{n}$ e $y_i\in\{0,1\}$ es \emph{linealmente separable} si existen $w\in\R^{n}$ y $b\in\R$ tales que $w\cdot x_i+b>0$ cuando $y_i=1$ y $w\cdot x_i+b<0$ cuando $y_i=0$.
\end{definicion}

\section{La regla de aprendizaje de Rosenblatt}

\begin{definicion}[Regla del perceptrón]\label{def:regla}
Sea $\eta>0$ una \emph{tasa de aprendizaje}. Se parte de $w=0$, $b=0$ y se recorren los ejemplos de forma cíclica. Para cada ejemplo $(x_i,y_i)$ se calcula $\hat y_i$ y se actualiza
\[
w\leftarrow w+\eta\,(y_i-\hat y_i)\,x_i,\qquad b\leftarrow b+\eta\,(y_i-\hat y_i).
\]
El algoritmo termina cuando una pasada completa no produce ningún error.
\end{definicion}

Como $y_i-\hat y_i\in\{-1,0,1\}$, los pesos solo cambian cuando el ejemplo está mal clasificado, y el cambio desplaza la preactivación en la dirección correcta: si $y_i=1$ y $\hat y_i=0$, la nueva preactivación de $x_i$ aumenta en $\eta(\norm{x_i}^{2}+1)>0$.

\addtocounter{definicion}{1}%

Para enunciar el teorema de convergencia conviene absorber el sesgo y usar etiquetas simétricas. Sea $\tilde x_i=(x_i,1)\in\R^{n+1}$, $\theta=(w,b)$ y $s_i=2y_i-1\in\{-1,+1\}$. El ejemplo $i$ está bien clasificado si $s_i\,\theta\cdot\tilde x_i>0$, y la regla de la \cref{def:regla} se escribe $\theta\leftarrow\theta+\eta\,s_i\,\tilde x_i$ cada vez que se comete un error. En particular, toda actualización se produce en un ejemplo con $s_i\,\theta\cdot\tilde x_i\le 0$.

\begin{teorema}[Convergencia del perceptrón; Novikoff]\label{teo:novikoff}
Supóngase que existen $u\in\R^{n+1}$ con $\norm{u}=1$ y $\gamma>0$ tales que $s_i\,u\cdot\tilde x_i\ge\gamma$ para todo $i$, y sea $R=\max_i\norm{\tilde x_i}$. Entonces, cualquiera que sea el orden en que se presentan los ejemplos, el número total $M$ de actualizaciones de la regla del perceptrón cumple
\[
M\le\frac{R^{2}}{\gamma^{2}}.
\]
\end{teorema}

\begin{observacion}
La cota no depende de $\eta$ ni del número de ejemplos, sino solo de la geometría de los datos: el cociente $R/\gamma$ mide lo «estrecho» que es el pasillo que separa las dos clases. Si los datos no son linealmente separables, no existe tal $u$ y la regla no termina nunca.
\end{observacion}

\section{El problema XOR}

La función XOR vale~1 en $(0,1)$ y $(1,0)$ y vale~0 en $(0,0)$ y $(1,1)$.

\begin{proposicion}\label{prop:xor}
El conjunto $\{((0,0),0),((0,1),1),((1,0),1),((1,1),0)\}$ no es linealmente separable. En consecuencia, ningún perceptrón calcula XOR.
\end{proposicion}

Minsky y Papert~\cite{minsky} analizaron en 1969 esta y otras limitaciones de los perceptrones de una capa, y su libro contribuyó a frenar la investigación en redes neuronales durante más de una década~\cite[cap.~1]{gbc}.

\section{Una capa oculta resuelve XOR}

La salida consiste en componer dos capas. Sea $\ReLU(t)=\max(0,t)$ y considérese la red de la \cref{fig:xor}, tomada de~\cite[§6.1]{gbc}:
\begin{equation}\label{eq:xor}
h_1=\ReLU(x_1+x_2),\qquad h_2=\ReLU(x_1+x_2-1),\qquad \hat y=h_1-2h_2.
\end{equation}

\addtocounter{figure}{1}%

En el espacio de la capa oculta, las cuatro entradas se transforman en $(0,0)$, $(1,0)$, $(1,0)$ y $(2,1)$: las dos entradas con salida~1 caen en el mismo punto, y la función lineal $h_1-2h_2$ separa las clases. La capa oculta no clasifica: cambia la \emph{representación} para que una función lineal pueda hacerlo. Aprender la representación a la vez que el clasificador es la idea que distingue el aprendizaje profundo de los modelos lineales del módulo~8.

\section{El perceptrón multicapa}

\begin{definicion}[Perceptrón multicapa]\label{def:mlp}
Sean $n_0,n_1,\dots,n_L$ enteros positivos. Un \emph{perceptrón multicapa} (MLP) de $L$ capas con anchuras $n_0\to n_1\to\dots\to n_L$ está formado por matrices $W_\ell\in\R^{n_{\ell-1}\times n_\ell}$, sesgos $b_\ell\in\R^{1\times n_\ell}$ y activaciones $\phi_\ell$, $\ell=1,\dots,L$. Para $x\in\R^{1\times n_0}$ se define $h^{(0)}=x$ y
\[
z^{(\ell)}=h^{(\ell-1)}W_\ell+b_\ell,\qquad h^{(\ell)}=\phi_\ell\bigl(z^{(\ell)}\bigr),\qquad \ell=1,\dots,L,
\]
con $\phi_\ell$ aplicada componente a componente. Las capas $1,\dots,L-1$ se llaman \emph{ocultas}. En un clasificador de $K$ clases, $n_L=K$, $\phi_L$ es la identidad y la salida $z^{(L)}$ es el vector de logits, que se convierte en $p=\softmax(z^{(L)})$.
\end{definicion}

\addtocounter{figure}{1}%

Con un lote de $B$ ejemplos apilados como filas de una matriz $X\in\R^{B\times n_0}$, la misma fórmula $XW_\ell+b_\ell$ (con $b_\ell$ sumado a cada fila) calcula todas las preactivaciones de la capa con un único producto de matrices. Esa es la razón práctica de la convención de vectores fila.

\begin{proposicion}[Número de parámetros]\label{prop:parametros-mlp}
Un MLP de anchuras $n_0\to n_1\to\dots\to n_L$ tiene $\sum_{\ell=1}^{L}(n_{\ell-1}+1)\,n_\ell$ parámetros.
\end{proposicion}

\addtocounter{definicion}{1}%

\begin{proposicion}[Colapso de las capas lineales]\label{prop:colapso}
Si todas las activaciones de un MLP son la identidad, existen $W\in\R^{n_0\times n_L}$ y $b\in\R^{1\times n_L}$ tales que $z^{(L)}=xW+b$ para toda entrada $x$. Además, $\operatorname{rg}W\le\min_{\ell}n_\ell$.
\end{proposicion}

La \cref{prop:colapso} explica por qué la no linealidad es imprescindible: sin ella, añadir capas no aumenta lo que la red puede representar, y una capa intermedia estrecha incluso lo reduce.

\addtocounter{definicion}{1}%


% =====================================================================
\chapter{Funciones de activación}
\label{tema:2}
% =====================================================================

La activación determina qué funciones puede representar una red y cómo fluye el gradiente a través de ella. Este tema estudia las activaciones clásicas (sigmoide y tangente hiperbólica), la ReLU, que se impuso en la década de 2010, y las activaciones suaves y con puerta que usan los modelos de lenguaje actuales. Se sigue la §6.3 de~\cite{gbc}, la §5.1.2 de~\cite{d2l} y las notas «Neural Networks Part~1» de~\cite{cs231n}; las definiciones de GELU y SwiGLU proceden de los artículos originales de Hendrycks y Gimpel~\cite{hendrycks} y de Shazeer~\cite{shazeer}.

\section{Sigmoide y tangente hiperbólica}

\begin{definicion}
La \emph{sigmoide logística} y la \emph{tangente hiperbólica} son
\[
\sigma(t)=\frac{1}{1+e^{-t}},\qquad \tanh(t)=\frac{e^{t}-e^{-t}}{e^{t}+e^{-t}}.
\]
\end{definicion}

\begin{proposicion}\label{prop:sigmoide}
Para todo $t\in\R$:
\begin{enumerate}
  \item $\sigma'(t)=\sigma(t)\bigl(1-\sigma(t)\bigr)$ y $0<\sigma'(t)\le 1/4$, con igualdad solo en $t=0$;
  \item $\tanh(t)=2\sigma(2t)-1$;
  \item $\tanh'(t)=1-\tanh^{2}(t)\in(0,1]$.
\end{enumerate}
\end{proposicion}

Por (ii), una red con activaciones tanh y otra con sigmoides representan la misma familia de funciones: el factor~2 y el desplazamiento $-1$ se absorben en los pesos y sesgos de las capas adyacentes (\cref{eje:2-tanh}). La diferencia práctica es que tanh está centrada en~0.

\begin{observacion}[Saturación y desvanecimiento del gradiente]\label{obs:desvanecimiento}
Cuando $\abs{t}$ es grande, $\sigma'(t)$ y $\tanh'(t)$ son casi nulas: la unidad está \emph{saturada}. En la retropropagación (\cref{tema:5}) el gradiente que atraviesa una capa se multiplica por $\phi'(z)$; con sigmoides, tras $k$ capas el factor es a lo sumo $4^{-k}$, que para $k=10$ vale menos de $10^{-6}$. Este \emph{desvanecimiento del gradiente} dificultó durante años el entrenamiento de redes profundas~\cite[§6.3.2]{gbc}.
\end{observacion}

\section{La unidad lineal rectificada}

\begin{definicion}[ReLU]
La \emph{unidad lineal rectificada} es $\ReLU(t)=\max(0,t)$. Su derivada es $\ReLU'(t)=\ind[t>0]$ para $t\neq 0$; en $t=0$ se toma por convenio el valor~0.
\end{definicion}

La ReLU tiene tres ventajas~\cite[§6.3.1]{gbc}: su cálculo es una comparación; en la región activa su derivada vale exactamente~1, de modo que el gradiente no se encoge al atravesar muchas capas; y es \emph{positivamente homogénea}, $\ReLU(ct)=c\ReLU(t)$ para $c\ge 0$, lo que simplifica el análisis de la varianza (\cref{tema:4}). Su inconveniente es que una unidad cuya preactivación es negativa para todos los datos tiene gradiente nulo y deja de aprender: se dice que está \emph{muerta}~\cite{cs231n}. Las variantes \emph{leaky} ReLU, $\max(\alpha t,t)$ con $0<\alpha<1$ fijo, y PReLU, con $\alpha$ aprendido~\cite{he}, evitan ese problema.

\section{Activaciones suaves: GELU y SiLU}

\begin{definicion}[GELU y SiLU]\label{def:gelu}
Sea $\Phi$ la función de distribución de la normal estándar. Se definen
\[
\GELU(t)=t\,\Phi(t),\qquad \SiLU(t)=t\,\sigma(t).
\]
Más en general, $\mathrm{Swish}_\beta(t)=t\,\sigma(\beta t)$ con $\beta>0$, de modo que $\SiLU=\mathrm{Swish}_1$.
\end{definicion}

Hendrycks y Gimpel~\cite{hendrycks} motivan GELU como el valor esperado de una puerta aleatoria: si $m\sim\mathrm{Bernoulli}(\Phi(t))$, entonces $\E[m\,t]=t\,\Phi(t)$. La entrada se deja pasar con una probabilidad que crece con su valor, en lugar de cortarse según su signo como en la ReLU. En la práctica se usa la aproximación
\[
\GELU(t)\approx\tfrac12\,t\Bigl(1+\tanh\Bigl[\sqrt{2/\pi}\,\bigl(t+0{,}044715\,t^{3}\bigr)\Bigr]\Bigr),
\]
cuyo error absoluto es menor que $5\cdot 10^{-4}$ en todo $\R$. GELU es la activación de BERT y de GPT-2.

\addtocounter{figure}{1}%

\begin{table}[h]
\centering\small
\caption{Valores de ReLU, GELU y SiLU.}\label{tab:activaciones}
\begin{tabular}{@{}lccccccc@{}}
\toprule
$t$ & $-2$ & $-1$ & $-0{,}5$ & $0$ & $0{,}5$ & $1$ & $2$ \\
\midrule
$\ReLU(t)$ & $0$ & $0$ & $0$ & $0$ & $0{,}5$ & $1$ & $2$ \\
$\GELU(t)$ & $-0{,}0455$ & $-0{,}1587$ & $-0{,}1543$ & $0$ & $0{,}3457$ & $0{,}8413$ & $1{,}9545$ \\
$\SiLU(t)$ & $-0{,}2384$ & $-0{,}2689$ & $-0{,}1888$ & $0$ & $0{,}3112$ & $0{,}7311$ & $1{,}7616$ \\
\bottomrule
\end{tabular}
\end{table}

\begin{proposicion}\label{prop:silu}
\begin{enumerate}
  \item $\SiLU'(t)=\sigma(t)\bigl(1+t\,(1-\sigma(t))\bigr)$; en particular $\SiLU'(0)=1/2$.
  \item $\GELU'(t)=\Phi(t)+t\,\varphi(t)$, donde $\varphi$ es la densidad normal estándar; en particular $\GELU'(0)=1/2$.
  \item Ni GELU ni SiLU son monótonas: ambas tienen un único mínimo global negativo, que vale aproximadamente $-0{,}170$ en $t\approx-0{,}752$ para GELU y $-0{,}278$ en $t\approx-1{,}278$ para SiLU.
\end{enumerate}
\end{proposicion}

Cuando $\beta\to\infty$, $\sigma(\beta t)\to\ind[t>0]$ para $t\neq 0$ y $\mathrm{Swish}_\beta$ tiende a la ReLU; el \cref{eje:2-swish-relu} prueba que la convergencia es uniforme.

\section{Unidades con puerta: GLU y SwiGLU}

En una unidad con puerta, una proyección de la entrada decide, componente a componente, cuánto se deja pasar de otra proyección.

\begin{definicion}[Capas prealimentadas con puerta; Shazeer]\label{def:swiglu}
Sean $d,H\ge 1$, $W_g,W_u\in\R^{d\times H}$ y $W_d\in\R^{H\times d}$. Para $x\in\R^{1\times d}$ se definen, sin sesgos,
\[
\FFN_{\mathrm{GLU}}(x)=\bigl(\sigma(xW_g)\odot xW_u\bigr)W_d,\qquad
\FFN_{\mathrm{SwiGLU}}(x)=\bigl(\SiLU(xW_g)\odot xW_u\bigr)W_d.
\]
La rama $xW_g$ es la \emph{puerta}, la rama $xW_u$ es el \emph{contenido} y $W_d$ proyecta de vuelta a dimensión $d$. El MLP clásico de un Transformer es $\FFN(x)=\phi(xW_1)W_2$, con $W_1\in\R^{d\times 4d}$ y $W_2\in\R^{4d\times d}$.
\end{definicion}

La unidad GLU fue propuesta por Dauphin \emph{et al.} para modelos de lenguaje convolucionales; Shazeer~\cite[§2]{shazeer} estudió variantes que sustituyen la sigmoide por otras funciones y observó que SwiGLU y GEGLU mejoran la calidad del Transformer frente a ReLU y GELU con el mismo número de parámetros.

\addtocounter{definicion}{1}%

\begin{proposicion}[Anchura de SwiGLU con igual número de parámetros]\label{prop:anchura}
Una capa $\FFN_{\mathrm{SwiGLU}}$ de anchura $H$ tiene $3dH$ pesos y un MLP clásico $d\to 4d\to d$ tiene $8d^{2}$. Ambos coinciden si y solo si $H=\tfrac{2}{3}\cdot 4d=8d/3$.
\end{proposicion}

Shazeer reduce la anchura en un factor $2/3$ por este motivo~\cite[§2]{shazeer}, y LLaMA adopta la misma regla: usa SwiGLU con dimensión oculta $\tfrac23\,4d$ en lugar de $4d$~\cite[§2.2]{touvron}. Para LLaMA-7B, con $d=4096$, se obtiene $8\cdot 4096/3\approx 10\,923$; el código publicado redondea al múltiplo de~256 siguiente, $H=11\,008$, con $3\cdot 4096\cdot 11\,008=135\,266\,304$ pesos por capa frente a $2\cdot 4096\cdot 16\,384=134\,217\,728$ del MLP clásico. Qué activación usan internamente Claude, Gemini o Grok no es información pública.


% =====================================================================
\chapter{Aproximación universal y profundidad}
\label{tema:3}
% =====================================================================

¿Qué funciones puede representar un perceptrón multicapa? Este tema demuestra que una sola capa oculta basta, en principio, para aproximar cualquier función continua en un compacto, y explica por qué ese resultado no justifica por sí solo el uso de redes anchas y poco profundas. La demostración en una dimensión es constructiva y usa solo la continuidad uniforme; el caso general se enuncia según Cybenko~\cite{cybenko}, Hornik~\cite{hornik} y Leshno \emph{et al.}~\cite{leshno}. La discusión sobre la profundidad sigue la §6.4.1 de~\cite{gbc} y el artículo de Telgarsky~\cite{telgarsky}; la §5.1.1.4 de~\cite{d2l} y el apartado «Representational power» de~\cite{cs231n} ofrecen el mismo punto de vista.

\section{Redes ReLU y funciones lineales a trozos}

\begin{definicion}[Función lineal a trozos]
Una función continua $f\colon\R\to\R$ es \emph{lineal a trozos} con \emph{codos} $t_1<\dots<t_m$ si es afín en cada uno de los intervalos $(-\infty,t_1]$, $[t_1,t_2]$, \dots, $[t_m,\infty)$. Cada uno de esos $m+1$ intervalos es una \emph{pieza}.
\end{definicion}

\begin{proposicion}\label{prop:piezas}
Sea $f(t)=c+\sum_{j=1}^{H}a_j\ReLU(w_jt+b_j)$ una red de una capa oculta con $H$ unidades ReLU, entrada y salida escalares. Entonces $f$ es continua y lineal a trozos con a lo sumo $H$ codos, es decir, con a lo sumo $H+1$ piezas.
\end{proposicion}

\begin{lema}[Sombrero]\label{lem:sombrero}
La función $\Lambda(t)=\ReLU(t)-2\ReLU(t-1)+\ReLU(t-2)$ vale $0$ para $t\le 0$ y para $t\ge 2$, vale $t$ en $[0,1]$ y $2-t$ en $[1,2]$.
\end{lema}

Sumando sombreros desplazados y escalados se obtiene cualquier poligonal con nodos equiespaciados; el teorema siguiente usa una construcción equivalente y más económica.

\section{El teorema en una dimensión}

\begin{teorema}[Aproximación universal en $\R$ con ReLU]\label{teo:au1d}
Sea $f\colon[a,b]\to\R$ continua y sea $\varepsilon>0$. Existe una red $g(t)=c+\sum_{j=1}^{H}a_j\ReLU(t-\tau_j)$ con una capa oculta tal que $\abs{f(t)-g(t)}<\varepsilon$ para todo $t\in[a,b]$.
\end{teorema}

\addtocounter{definicion}{1}%

\section{El teorema general}

\begin{teorema}[Aproximación universal; Cybenko, Hornik, Leshno \emph{et al.}]\label{teo:au}
Sea $\phi\colon\R\to\R$ continua. Las funciones de la forma
\[
g(x)=\sum_{j=1}^{H}a_j\,\phi\bigl(x\cdot w_j+b_j\bigr),\qquad H\in\mathbb{N},\ a_j,b_j\in\R,\ w_j\in\R^{n},
\]
son densas en el espacio de funciones continuas sobre cualquier compacto $K\subset\R^{n}$ con la norma del supremo si y solo si $\phi$ no es un polinomio.
\end{teorema}

Cybenko~\cite{cybenko} lo demostró en 1989 para activaciones sigmoidales usando el teorema de Hahn--Banach y el de representación de Riesz; Hornik~\cite{hornik} lo extendió a activaciones acotadas y no constantes, y Leshno, Lin, Pinkus y Schocken~\cite{leshno} obtuvieron la caracterización anterior. La demostración del caso general excede el nivel de este curso. La necesidad de la condición sí es elemental:

\begin{proposicion}\label{prop:polinomio}
Si $\phi$ es un polinomio de grado $k$, toda red de la forma del \cref{teo:au} es un polinomio de grado a lo sumo $k$, y las redes no son densas en $C([0,1])$.
\end{proposicion}

\section{Qué no dice el teorema: el papel de la profundidad}

El \cref{teo:au} es un resultado de existencia. No acota el número $H$ de unidades necesarias, que para algunas funciones crece exponencialmente con la dimensión $n$ de la entrada~\cite[§6.4.1]{gbc}; no garantiza que el descenso de gradiente encuentre los pesos, y no dice nada sobre la generalización a datos no vistos. La profundidad aporta una ventaja que el teorema no recoge: componer capas permite reutilizar piezas, y algunas funciones que una red profunda representa con pocas unidades requieren exponencialmente muchas en una red poco profunda. El ejemplo siguiente lo muestra de forma exacta.

\begin{proposicion}[Diente de sierra]\label{prop:sierra}
Sea $g(t)=2\ReLU(t)-4\ReLU(t-\tfrac12)$, que en $[0,1]$ vale $2t$ en $[0,\tfrac12]$ y $2-2t$ en $[\tfrac12,1]$, y sea $g^{(k)}=g\circ\cdots\circ g$ ($k$ veces). Entonces:
\begin{enumerate}
  \item en $[0,1]$, $g^{(k)}$ es lineal a trozos con exactamente $2^{k}$ piezas de longitud $2^{-k}$, y en cada una recorre todo $[0,1]$;
  \item $g^{(k)}$ se calcula con una red de $k$ capas ocultas de 2 unidades ReLU cada una;
  \item toda red de una capa oculta que coincida con $g^{(k)}$ en $[0,1]$ necesita al menos $2^{k}-1$ unidades.
\end{enumerate}
\end{proposicion}

Con $k=20$, la red profunda usa 40 unidades y la poco profunda más de un millón. Telgarsky~\cite{telgarsky} demostró que la separación persiste aunque solo se pida \emph{aproximar}: hay funciones calculables por redes de profundidad del orden de $k^{3}$ y anchura constante que una red de profundidad $O(k)$ solo aproxima con del orden de $2^{k}$ unidades. Los LLM son profundos por esta razón: LLaMA-7B apila 32 bloques~\cite[§2]{touvron}, cada uno con un MLP de anchura moderada.


% =====================================================================
\chapter{Inicialización y propagación de la varianza}
\label{tema:4}
% =====================================================================

El descenso de gradiente necesita un punto de partida. Si los pesos iniciales son demasiado grandes, las activaciones crecen capa a capa o saturan; si son demasiado pequeños, la señal se apaga y los gradientes con ella. Este tema muestra primero por qué los pesos no pueden inicializarse todos iguales y deduce después, a partir de la varianza de una suma de variables independientes (módulo~5), las reglas de inicialización de LeCun, Xavier y He. Se sigue la §8.4 de~\cite{gbc}, la §5.4 de~\cite{d2l} y el apartado «Weight Initialization» de las notas «Neural Networks Part~2» de~\cite{cs231n}; los resultados proceden de Glorot y Bengio~\cite{glorot} y de He \emph{et al.}~\cite{he}.

\section{La ruptura de la simetría}

\begin{proposicion}[Simetría de las unidades ocultas]\label{prop:simetria}
Sea un MLP con una capa oculta, $z=xW_1+b_1$, $h=\phi(z)$, $o=hW_2+b_2$, y sea $L$ cualquier pérdida que dependa de $o$. Si en la inicialización las columnas $i$ y $j$ de $W_1$ son iguales, $(b_1)_i=(b_1)_j$, y las filas $i$ y $j$ de $W_2$ son iguales, entonces tras cualquier número de pasos de descenso de gradiente (por lotes o estocástico) esas igualdades se mantienen.
\end{proposicion}

En particular, si todos los pesos se inicializan a~0 (o a una misma constante), todas las unidades ocultas calculan siempre lo mismo y la capa equivale a una sola unidad~\cite{cs231n}, \cite[§5.4.1.3]{d2l}. Por eso los pesos se inicializan con valores \emph{aleatorios}. Los sesgos, en cambio, pueden empezar en~0. En un modelo lineal, como la regresión logística, el problema no se plantea (\cref{eje:4-simetria}).

\section{La varianza de una preactivación}

\begin{lema}\label{lem:var-producto}
Sean $w$ y $h$ variables aleatorias independientes con $\E[w]=0$. Entonces $\E[wh]=0$ y $\Var(wh)=\Var(w)\,\E[h^{2}]$.
\end{lema}

\begin{proposicion}[Propagación hacia delante]\label{prop:var-adelante}
Sea $z=\sum_{i=1}^{\nin}w_ih_i$ la preactivación de una unidad. Supóngase que los $w_i$ son independientes entre sí y de los $h_i$, con media~0 y varianza $s^{2}$, y que los $h_i$ tienen el mismo segundo momento $\E[h^{2}]$. Entonces $\E[z]=0$ y
\[
\Var(z)=\nin\,s^{2}\,\E[h^{2}].
\]
\end{proposicion}

El objetivo es elegir $s^{2}$ para que $\Var(z)$ sea la misma en todas las capas. Lo que relaciona $\E[h^{2}]$ con la varianza de la capa anterior es la activación.

\section{Activaciones lineales o tanh: LeCun y Xavier}

Si la activación es la identidad, o es tanh y opera cerca de~0, donde $\tanh t\approx t$, entonces $\E[h^{2}]\approx\Var(z_{\text{anterior}})$ y la \cref{prop:var-adelante} da $\Var(z)=\nin s^{2}\Var(z_{\text{anterior}})$. Conservar la varianza exige
\[
s^{2}=\frac{1}{\nin},
\]
que es la inicialización de \emph{LeCun}, la misma calibración $1/\sqrt{n}$ que recomiendan las notas de~\cite{cs231n}.

\begin{proposicion}[Propagación hacia atrás]\label{prop:var-atras}
En la retropropagación (\cref{tema:5}), el gradiente respecto a la activación $h_i$ de una capa es $g_i=\sum_{j=1}^{\nout}W_{ij}\delta_j$. Con las mismas hipótesis de independencia y media nula, $\Var(g_i)=\nout\,s^{2}\,\E[\delta^{2}]$.
\end{proposicion}

Para que la varianza de los gradientes tampoco cambie de capa en capa hace falta $\nout s^{2}=1$. Las dos condiciones solo son compatibles si $\nin=\nout$; Glorot y Bengio~\cite{glorot} proponen el compromiso de tomar $1/s^{2}$ igual a la media aritmética de $\nin$ y $\nout$.

\begin{definicion}[Inicialización de Xavier o Glorot]\label{def:xavier}
Los pesos de una capa $\nin\to\nout$ se toman independientes con media~0 y varianza
\[
s^{2}=\frac{2}{\nin+\nout},
\]
bien con distribución normal $\mathcal{N}(0,s^{2})$, bien uniforme en $[-a,a]$ con $a=\sqrt{6/(\nin+\nout)}$.
\end{definicion}

\begin{lema}\label{lem:uniforme}
Si $w\sim U(-a,a)$, entonces $\Var(w)=a^{2}/3$.
\end{lema}

Con $a^{2}=6/(\nin+\nout)$, el \cref{lem:uniforme} da $a^{2}/3=2/(\nin+\nout)$, de modo que las dos versiones de la \cref{def:xavier} tienen la misma varianza.

\addtocounter{definicion}{1}%

\section{Activaciones ReLU: la inicialización de He}

\begin{lema}\label{lem:relu2}
Si $z$ es una variable aleatoria con densidad simétrica respecto a~0 y varianza $\sigma^{2}$, entonces $\E[\ReLU(z)^{2}]=\sigma^{2}/2$.
\end{lema}

\begin{teorema}[Inicialización de He]\label{teo:he}
En un MLP con activaciones ReLU, supóngase que las preactivaciones de cada capa tienen densidad simétrica y que se cumplen las hipótesis de la \cref{prop:var-adelante}. Entonces $\Var(z^{(\ell+1)})=\tfrac12\,\nin s^{2}\,\Var(z^{(\ell)})$, y la varianza se conserva capa a capa si y solo si
\[
s^{2}=\frac{2}{\nin}.
\]
\end{teorema}

He \emph{et al.}~\cite{he} obtienen esta regla y la generalizan a PReLU; las notas de~\cite{cs231n} la presentan como la recomendación práctica para redes ReLU. El factor~2 compensa exactamente que la ReLU anula la mitad del segundo momento. Con la inicialización de LeCun o de Xavier ($s^{2}=1/n$ cuando $\nin=\nout=n$), en cambio, la varianza se divide por~2 en cada capa ReLU y la desviación típica por $\sqrt2$.

\addtocounter{definicion}{1}%

\section{Un experimento}

El \cref{tab:experimento} muestra la desviación típica de las activaciones de una red de diez capas de anchura~256, con entradas normales estándar, para cinco inicializaciones (media de 20 repeticiones con 1000 entradas cada una).

\begin{table}[h]
\centering\small
\caption{Desviación típica de las activaciones en una red de diez capas de anchura~256.}\label{tab:experimento}
\begin{tabular}{@{}llcccl@{}}
\toprule
Pesos & Activación & Capa 1 & Capa 4 & Capa 10 & Diagnóstico \\
\midrule
$\mathcal{N}(0,1)$ & tanh & $0{,}975$ & $0{,}974$ & $0{,}974$ & saturada: $86{,}5\,\%$ con $\abs{h}>0{,}99$ \\
$\mathcal{N}(0;\,0{,}01^{2})$ & tanh & $0{,}156$ & $6{,}4\cdot 10^{-4}$ & $1{,}1\cdot 10^{-8}$ & señal apagada \\
Xavier & tanh & $0{,}628$ & $0{,}357$ & $0{,}228$ & decae despacio \\
Xavier & ReLU & $0{,}584$ & $0{,}205$ & $0{,}0254$ & se divide por $\sqrt2$ en cada capa \\
He & ReLU & $0{,}826$ & $0{,}819$ & $0{,}812$ & estable \\
\bottomrule
\end{tabular}
\end{table}

Los números concuerdan con la teoría. Con He, la preactivación de la capa~1 es $\mathcal{N}(0,2)$ y la desviación típica de $\ReLU(z)$ es $\sqrt{1-1/\pi}=0{,}826$. Con Xavier y ReLU, el factor $2^{-1/2}$ por capa predice $0{,}584\cdot 2^{-9/2}=0{,}0258$ en la capa~10. Con $\mathcal{N}(0,1)$ y tanh la señal no desaparece, pero las unidades están saturadas: la derivada media $1-h^{2}$ vale $0{,}051$, y los gradientes se multiplican por números pequeños en cada capa (\cref{obs:desvanecimiento}).

\addtocounter{definicion}{1}%


% =====================================================================
\chapter{Retropropagación en un perceptrón multicapa}
\label{tema:5}
% =====================================================================

Entrenar una red por descenso de gradiente (módulo~7) exige calcular el gradiente de la pérdida respecto a todos los parámetros. La \emph{retropropagación} lo hace aplicando la regla de la cadena del módulo~4 desde la salida hacia la entrada y reutilizando cada derivada intermedia. Este tema la deduce para un MLP con salida softmax y entropía cruzada y la aplica a mano a una red pequeña. Se sigue la §6.5 de~\cite{gbc}, en particular la §6.5.4, las §6.6.2--6.6.4 de~\cite{jm3}, el capítulo~8 de~\cite{bishopdl}, la §5.3 de~\cite{d2l}, las notas «Backpropagation, Intuitions» de~\cite{cs231n} y la clase~3 de CS224N~\cite{cs224n}. El método se popularizó con el artículo de Rumelhart, Hinton y Williams~\cite{rumelhart}; su generalización a grafos arbitrarios, la diferenciación automática, es el objeto del módulo~10.

\section{Grafos computacionales}

\begin{definicion}[Grafo computacional]
Un \emph{grafo computacional} es un grafo dirigido acíclico cuyos nodos son variables (escalares, vectores o matrices) y en el que cada nodo no inicial se obtiene aplicando una operación conocida a sus predecesores. La \emph{pasada hacia delante} evalúa los nodos en un orden topológico; la \emph{pasada hacia atrás} recorre el grafo en sentido inverso.
\end{definicion}

\begin{proposicion}[Regla de la cadena en un grafo]\label{prop:cadena-grafo}
Sea $L$ un nodo escalar final y $v$ un nodo con sucesores $u_1,\dots,u_r$. Entonces
\[
\frac{\partial L}{\partial v}=\sum_{k=1}^{r}\frac{\partial L}{\partial u_k}\,\frac{\partial u_k}{\partial v}.
\]
\end{proposicion}

La consecuencia algorítmica es que, si se conocen $\partial L/\partial u_k$ para todos los sucesores de $v$, el gradiente respecto a $v$ se obtiene con operaciones locales. Recorriendo el grafo en orden topológico inverso, cada derivada se calcula una sola vez~\cite[§6.5.3]{gbc}. Cuando un nodo alimenta a varios sucesores, sus contribuciones se \emph{suman}.

\section{El gradiente de la softmax con entropía cruzada}

\begin{proposicion}\label{prop:delta-salida}
Sea $z\in\R^{1\times K}$, $p=\softmax(z)$ e $y\in\{0,\dots,K-1\}$, y sea $L=-\ln p_y$. Entonces
\[
\frac{\partial L}{\partial z_k}=p_k-\ind[k=y],\qquad\text{es decir,}\qquad \nabla_{\!z}L=p-e_y.
\]
\end{proposicion}

\begin{corolario}\label{cor:suma-cero}
Las componentes de $\nabla_{\!z}L$ suman~0.
\end{corolario}

El gradiente tiene una lectura directa: el descenso sube el logit de la clase correcta, que es la única componente negativa, y baja los demás en proporción a la probabilidad que el modelo les asignaba por error. La misma expresión $p-e_y$ es la que se retropropaga en cada posición durante el preentrenamiento de un LLM.

\section{Capas lineales y de activación}

\begin{proposicion}[Gradientes de una capa afín]\label{prop:capa-lineal}
Sea $z=xW+b$ con $x\in\R^{1\times n}$, $W\in\R^{n\times m}$ y $b\in\R^{1\times m}$, y sea $g=\nabla_{\!z}L\in\R^{1\times m}$. Entonces
\[
\nabla_{\!W}L=x^{\top}g\in\R^{n\times m},\qquad \nabla_{\!b}L=g,\qquad \nabla_{\!x}L=g\,W^{\top}\in\R^{1\times n}.
\]
Para un lote $X\in\R^{B\times n}$ con gradientes $G\in\R^{B\times m}$ por filas y pérdida igual a la suma de las pérdidas, $\nabla_{\!W}L=X^{\top}G$, $\nabla_{\!b}L=\mathbf{1}^{\top}G$ y $\nabla_{\!X}L=GW^{\top}$.
\end{proposicion}

\begin{observacion}[Análisis dimensional]
El gradiente respecto a una matriz tiene siempre la forma de la matriz. En la \cref{prop:capa-lineal}, $x^{\top}$ es $n\times 1$ y $g$ es $1\times m$, de modo que $x^{\top}g$ es $n\times m$, como $W$. Comprobar las formas basta a menudo para recordar el orden y las trasposiciones: «entrada traspuesta por gradiente de la salida»~\cite{cs231n}. Con vectores columna, como en las notas de~\cite{cs231n}, las mismas fórmulas aparecen traspuestas.
\end{observacion}

\begin{proposicion}[Gradiente a través de una activación]\label{prop:activacion}
Sea $h=\phi(z)$ aplicada componente a componente. Entonces $\nabla_{\!z}L=\nabla_{\!h}L\odot\phi'(z)$.
\end{proposicion}

\section{Las ecuaciones de la retropropagación}

\begin{teorema}[Retropropagación en un MLP de dos capas]\label{teo:backprop}
Sea el MLP $z^{(1)}=xW_1+b_1$, $h=\phi(z^{(1)})$, $z^{(2)}=hW_2+b_2$, $p=\softmax(z^{(2)})$, $L=-\ln p_y$. Definiendo $\delta^{(2)}=\nabla_{z^{(2)}}L$ y $\delta^{(1)}=\nabla_{z^{(1)}}L$, se cumple
\[
\begin{aligned}
\delta^{(2)}&=p-e_y, & \nabla_{\!W_2}L&=h^{\top}\delta^{(2)}, & \nabla_{\!b_2}L&=\delta^{(2)},\\
\nabla_{\!h}L&=\delta^{(2)}W_2^{\top}, & \delta^{(1)}&=\nabla_{\!h}L\odot\phi'(z^{(1)}), & & \\
\nabla_{\!W_1}L&=x^{\top}\delta^{(1)}, & \nabla_{\!b_1}L&=\delta^{(1)}, & \nabla_{\!x}L&=\delta^{(1)}W_1^{\top}.
\end{aligned}
\]
\end{teorema}

\addtocounter{figure}{1}%

\begin{observacion}[Coste]
La pasada hacia delante realiza dos productos de matrices; la hacia atrás, tres para los gradientes de los parámetros y $\nabla_hL$, más uno opcional para $\nabla_xL$, todos de las mismas dimensiones. El coste de calcular el gradiente completo es, por tanto, del mismo orden que el de evaluar la red: aproximadamente el doble. Calcular cada derivada parcial por separado, con diferencias finitas, costaría una evaluación por parámetro.
\end{observacion}

\section{Un ejemplo completo}

\addtocounter{definicion}{1}%

\section{Comprobación numérica del gradiente}

Un error en una sola ecuación de la retropropagación produce gradientes incorrectos sin ningún mensaje de error. La comprobación estándar compara con diferencias finitas~\cite{cs231n}.

\begin{proposicion}[Diferencias centradas]\label{prop:centradas}
Sea $f\colon\R\to\R$ de clase $C^{3}$ en un entorno de $\theta$. Entonces
\[
\frac{f(\theta+\varepsilon)-f(\theta-\varepsilon)}{2\varepsilon}=f'(\theta)+O(\varepsilon^{2}).
\]
\end{proposicion}

Aplicada a cada entrada de $W_1$ en el \cref{ej:backprop} con $\varepsilon=10^{-6}$, la fórmula reproduce $\nabla_{\!W_1}L$ con un error menor que $10^{-9}$. La función ReLU no es derivable en~0; la comprobación es fiable mientras ninguna preactivación esté a distancia menor que $\varepsilon$ de ese punto.


% =====================================================================
\chapter{El modelo de lenguaje neuronal}
\label{tema:6}
% =====================================================================

Este tema reúne los anteriores en un modelo de lenguaje. Bengio, Ducharme, Vincent y Jauvin~\cite{bengio} propusieron en 2003 sustituir la tabla de un modelo de $n$-gramas por una red: cada token del contexto se convierte en un vector, los vectores se concatenan y un MLP produce la distribución del token siguiente. La presentación sigue la §2 del artículo original y la §6.5 de~\cite{jm3}; la motivación por el tamaño de las tablas, la §9.3 de~\cite{d2l}. Los resultados numéricos son los de la práctica del módulo, con el corpus de nombres del módulo~1.

\section{La maldición de la dimensionalidad}

Un modelo de $n$-gramas sobre un vocabulario de $K$ tokens tiene $K^{n}$ entradas (módulo~1, tema~2). Con $K=28$ y un contexto de tres caracteres ($n=4$), la tabla tiene $28^{4}=614\,656$ casillas, frente a los $1885$ ejemplos de entrenamiento del corpus de nombres: casi todos los contextos no se han visto nunca y su estimación depende por completo del suavizado. El problema de fondo es que la tabla trata cada contexto como un símbolo aislado. Los contextos «.an» y «.en» no comparten nada, aunque las vocales \texttt{a} y \texttt{e} aparezcan en entornos parecidos.

\section{Embeddings como capa lineal}

\begin{definicion}[Matriz de embeddings]
Sea $\V$ un vocabulario de $K$ tokens y $d\ge 1$. Una \emph{matriz de embeddings} es $C\in\R^{K\times d}$; el \emph{embedding} del token de índice $i$ es su fila $C[i]=e_i^{\top}C\in\R^{1\times d}$, donde $e_i$ es el vector \emph{one-hot} de $i$ (módulo~2, tema~5).
\end{definicion}

\begin{proposicion}[Gradiente de los embeddings]\label{prop:grad-embedding}
Sea un contexto de tokens $(x_1,\dots,x_c)$ y sea $e=(C[x_1],\dots,C[x_c])\in\R^{1\times cd}$ la concatenación de sus embeddings. Si $\nabla_{\!e}L=(g_1,\dots,g_c)$ con $g_j\in\R^{1\times d}$, entonces
\[
\nabla_{\!C[i]}L=\sum_{j\,:\,x_j=i}g_j\qquad\text{para cada token } i,
\]
y las filas de los tokens que no aparecen en el contexto tienen gradiente nulo.
\end{proposicion}

En la práctica, las bibliotecas no multiplican por vectores \emph{one-hot}: leen la fila correspondiente y, en la pasada hacia atrás, acumulan los gradientes en las filas usadas (en NumPy, con \texttt{np.add.at}).

\section{El modelo de Bengio \emph{et al.}}

\begin{definicion}[Modelo de lenguaje neuronal prealimentado]\label{def:bengio}
Sea $n\ge 2$ el orden del modelo y $c=n-1$ la longitud del contexto. Dado el contexto $(x_{t-c},\dots,x_{t-1})$, se calcula
\[
e=\bigl(C[x_{t-1}],\dots,C[x_{t-c}]\bigr)\in\R^{1\times cd},\qquad
y=b+eW+\tanh(\tilde b+eH)\,U\in\R^{1\times K},
\]
con $H\in\R^{cd\times h}$, $\tilde b\in\R^{1\times h}$, $U\in\R^{h\times K}$, $b\in\R^{1\times K}$ y una matriz opcional de \emph{conexiones directas} $W\in\R^{cd\times K}$, y se toma $P(x_t=k\mid\text{contexto})=\softmax(y)_k$~\cite[§2]{bengio}.
\end{definicion}

\begin{proposicion}[Número de parámetros]\label{prop:param-bengio}
Con conexiones directas, el modelo de la \cref{def:bengio} tiene
\[
K(1+nd+h)+h\bigl(1+(n-1)d\bigr)
\]
parámetros; sin ellas, $K(1+nd+h)+h(1+(n-1)d)-K(n-1)d$.
\end{proposicion}

\addtocounter{figure}{1}%

\addtocounter{definicion}{1}%

\section{Entrenamiento y regularización}

Los parámetros $\theta$ se ajustan minimizando la entropía cruzada media sobre los $N$ pares (contexto, token siguiente) de entrenamiento, con decaimiento de pesos (módulo~7):
\[
J(\theta)=-\frac1N\sum_{r=1}^{N}\ln P_\theta\bigl(x^{(r)}_t\mid\text{contexto}^{(r)}\bigr)+\frac{\lambda}{2}\norm{\theta}^{2},
\qquad \theta\leftarrow\theta-\eta\bigl(\nabla\mathcal{L}(\theta)+\lambda\theta\bigr),
\]
donde $\mathcal{L}$ es el primer término de $J$. El gradiente de $\mathcal{L}$ se obtiene con el \cref{teo:backprop}, aplicado al lote completo, y con la \cref{prop:grad-embedding} para $C$. La perplejidad es $\PPL=e^{\mathcal{L}}$ (módulo~1, tema~5).

\begin{table}[h]
\centering\small
\caption{Perplejidades en el corpus de nombres del módulo~1 (1000 pasos de descenso por lotes completos, $\eta=0{,}5$).}\label{tab:resultados}
\begin{tabular}{@{}lcrcc@{}}
\toprule
Modelo & Contexto & Parámetros & \multicolumn{2}{c@{}}{PPL} \\
\cmidrule(l){4-5}
 & & & entren. & valid. \\
\midrule
Uniforme & 0 & 0 & 28 & 28 \\
Bigramas, suavizado $\alpha=0{,}2$ & 1 & 784 & --- & $12{,}42$ \\
MLP, $d=10$, $H=64$, $\lambda=0{,}005$ & 3 & 4084 & $6{,}62$ & $10{,}45$ \\
MLP sin decaimiento ($\lambda=0$) & 3 & 4084 & $3{,}76$ & $30{,}73$ \\
Tabla de 4-gramas & 3 & 614\,656 & --- & (casi vacía) \\
\bottomrule
\end{tabular}
\end{table}

El MLP regularizado reduce la perplejidad de validación del modelo de bigramas en un $16\,\%$, de $12{,}42$ a $10{,}45$, con menos de una centésima parte de los parámetros de una tabla de 4-gramas. Sin decaimiento de pesos, la red memoriza los nombres de entrenamiento: la perplejidad de entrenamiento baja a $3{,}76$ y la de validación sube a $30{,}73$, peor que el modelo uniforme. Es el sobreajuste del módulo~8.

\section{Por qué generaliza}

La red no ve tokens, sino vectores. Si dos tokens tienen embeddings próximos, los contextos que difieren solo en ellos producen distribuciones próximas. El resultado siguiente lo cuantifica. Se usa la norma espectral $\norm{A}_2=\max_{\norm{v}=1}\norm{vA}$ (módulo~3).

\begin{proposicion}\label{prop:lipschitz}
Sea $f(e)=\ReLU(eW_1+b_1)W_2+b_2$ la función que lleva la concatenación de embeddings a los logits, y sea $p(e)=\softmax(f(e))$. Para todo par $e,e'$ y todo token $k$,
\[
\norm{f(e)-f(e')}\le\norm{W_1}_2\norm{W_2}_2\,\norm{e-e'},\qquad
\abs{\ln p_k(e)-\ln p_k(e')}\le 2\,\norm{W_1}_2\norm{W_2}_2\,\norm{e-e'}.
\]
\end{proposicion}

Los embeddings se aprenden junto con el resto de la red, y tokens que aparecen en contextos parecidos reciben gradientes parecidos (\cref{prop:grad-embedding}); en la práctica del módulo, las vocales acaban agrupadas sin que nadie haya indicado a la red qué es una vocal. Lo aprendido para un contexto se transfiere así, por la \cref{prop:lipschitz}, a todos los contextos cercanos.

\addtocounter{definicion}{1}%



\backmatter
% =====================================================================
\begin{thebibliography}{99}
\addcontentsline{toc}{chapter}{Bibliografía}
\markboth{Bibliografía}{}
\small

\bibitem{bengio}
Y.~Bengio, R.~Ducharme, P.~Vincent y C.~Jauvin.
\newblock A neural probabilistic language model.
\newblock \emph{Journal of Machine Learning Research}, 3:1137--1155, 2003.
\newblock \url{https://www.jmlr.org/papers/v3/bengio03a.html}

\bibitem{bishopdl}
C.~M. Bishop y H.~Bishop.
\newblock \emph{Deep Learning: Foundations and Concepts}.
\newblock Springer, 2024.
\newblock Cap.~6, «Deep Neural Networks»; cap.~7, «Gradient Descent»; cap.~8, «Backpropagation».
\newblock \url{https://www.bishopbook.com/}

\bibitem{cmu11785}
B.~Raj \emph{et al.}
\newblock 11-785: Introduction to Deep Learning.
\newblock Carnegie Mellon University, otoño de 2025.
\newblock \url{https://deeplearning.cs.cmu.edu/}

\bibitem{cybenko}
G.~Cybenko.
\newblock Approximation by superpositions of a sigmoidal function.
\newblock \emph{Mathematics of Control, Signals and Systems}, 2(4):303--314, 1989.

\bibitem{glorot}
X.~Glorot e Y.~Bengio.
\newblock Understanding the difficulty of training deep feedforward neural networks.
\newblock En \emph{Proceedings of the Thirteenth International Conference on Artificial Intelligence and Statistics (AISTATS)}, PMLR~9:249--256, 2010.
\newblock \url{https://proceedings.mlr.press/v9/glorot10a.html}

\bibitem{gbc}
I.~Goodfellow, Y.~Bengio y A.~Courville.
\newblock \emph{Deep Learning}.
\newblock MIT Press, 2016.
\newblock Cap.~6, «Deep Feedforward Networks», y §8.4, «Parameter Initialization Strategies».
\newblock \url{https://www.deeplearningbook.org/}

\bibitem{he}
K.~He, X.~Zhang, S.~Ren y J.~Sun.
\newblock Delving deep into rectifiers: Surpassing human-level performance on ImageNet classification.
\newblock En \emph{Proceedings of the IEEE International Conference on Computer Vision (ICCV)}, 2015.
\newblock \url{https://arxiv.org/abs/1502.01852}

\bibitem{hendrycks}
D.~Hendrycks y K.~Gimpel.
\newblock Gaussian error linear units (GELUs).
\newblock Prepublicación arXiv:1606.08415, 2016.
\newblock \url{https://arxiv.org/abs/1606.08415}

\bibitem{hornik}
K.~Hornik.
\newblock Approximation capabilities of multilayer feedforward networks.
\newblock \emph{Neural Networks}, 4(2):251--257, 1991.

\bibitem{jm3}
D.~Jurafsky y J.~H. Martin.
\newblock \emph{Speech and Language Processing}, 3.ª ed., borrador del 19 de agosto de 2026.
\newblock Cap.~6, «Neural Networks».
\newblock \url{https://web.stanford.edu/~jurafsky/slp3/}

\bibitem{leshno}
M.~Leshno, V.~Y. Lin, A.~Pinkus y S.~Schocken.
\newblock Multilayer feedforward networks with a nonpolynomial activation function can approximate any function.
\newblock \emph{Neural Networks}, 6(6):861--867, 1993.

\bibitem{minsky}
M.~Minsky y S.~Papert.
\newblock \emph{Perceptrons: An Introduction to Computational Geometry}.
\newblock MIT Press, 1969.

\bibitem{novikoff}
A.~B.~J. Novikoff.
\newblock On convergence proofs on perceptrons.
\newblock En \emph{Proceedings of the Symposium on the Mathematical Theory of Automata}, vol.~12, págs.~615--622. Polytechnic Institute of Brooklyn, 1962.

\bibitem{radford2018}
A.~Radford, K.~Narasimhan, T.~Salimans e I.~Sutskever.
\newblock Improving language understanding by generative pre-training.
\newblock Informe técnico, OpenAI, 2018.
\newblock \url{https://cdn.openai.com/research-covers/language-unsupervised/language_understanding_paper.pdf}

\bibitem{radford2019}
A.~Radford, J.~Wu, R.~Child, D.~Luan, D.~Amodei e I.~Sutskever.
\newblock Language models are unsupervised multitask learners.
\newblock Informe técnico, OpenAI, 2019.
\newblock \url{https://cdn.openai.com/better-language-models/language_models_are_unsupervised_multitask_learners.pdf}

\bibitem{rosenblatt}
F.~Rosenblatt.
\newblock The perceptron: A probabilistic model for information storage and organization in the brain.
\newblock \emph{Psychological Review}, 65(6):386--408, 1958.

\bibitem{rumelhart}
D.~E. Rumelhart, G.~E. Hinton y R.~J. Williams.
\newblock Learning representations by back-propagating errors.
\newblock \emph{Nature}, 323:533--536, 1986.

\bibitem{shazeer}
N.~Shazeer.
\newblock GLU variants improve Transformer.
\newblock Prepublicación arXiv:2002.05202, 2020.
\newblock \url{https://arxiv.org/abs/2002.05202}

\bibitem{mit18065}
G.~Strang.
\newblock 18.065: Matrix Methods in Data Analysis, Signal Processing, and Machine Learning.
\newblock Massachusetts Institute of Technology, MIT OpenCourseWare, primavera de 2018.
\newblock Lección~27, «Backpropagation»; problemas de la sección VII.2 del libro \emph{Linear Algebra and Learning from Data}.
\newblock \url{https://ocw.mit.edu/courses/18-065-matrix-methods-in-data-analysis-signal-processing-and-machine-learning-spring-2018/}

\bibitem{telgarsky}
M.~Telgarsky.
\newblock Benefits of depth in neural networks.
\newblock En \emph{Conference on Learning Theory (COLT)}, 2016.
\newblock \url{https://arxiv.org/abs/1602.04485}

\bibitem{touvron}
H.~Touvron, T.~Lavril, G.~Izacard, X.~Martinet \emph{et al.}
\newblock LLaMA: Open and efficient foundation language models.
\newblock Prepublicación arXiv:2302.13971, 2023.
\newblock \url{https://arxiv.org/abs/2302.13971}

\bibitem{vaswani}
A.~Vaswani, N.~Shazeer, N.~Parmar, J.~Uszkoreit, L.~Jones, A.~N. Gomez, Ł.~Kaiser e I.~Polosukhin.
\newblock Attention is all you need.
\newblock En \emph{Advances in Neural Information Processing Systems 30}, 2017.
\newblock \url{https://arxiv.org/abs/1706.03762}

\bibitem{cs4780}
K.~Q. Weinberger.
\newblock CS4780/5780: Machine Learning for Intelligent Systems.
\newblock Cornell University, otoño de 2018. Lección~3, «The Perceptron».
\newblock \url{https://www.cs.cornell.edu/courses/cs4780/2018fa/lectures/lecturenote03.html}

\bibitem{cs224n}
D.~Yang e Y.~Choi.
\newblock CS224N: Natural Language Processing with Deep Learning.
\newblock Stanford University, invierno de 2026. Clase~3, «Backpropagation and Neural Network Basics», y clase~4, «Language Models and RNNs».
\newblock \url{https://web.stanford.edu/class/cs224n/}

\bibitem{cs224n2018}
C.~Manning y R.~Socher.
\newblock CS224N: Natural Language Processing with Deep Learning.
\newblock Stanford University, invierno de 2018. Tarea~1.
\newblock \url{https://web.stanford.edu/class/archive/cs/cs224n/cs224n.1184/assignment1/assignment1.pdf}

\bibitem{yang}
Z.~Yang, Z.~Dai, R.~Salakhutdinov y W.~W. Cohen.
\newblock Breaking the softmax bottleneck: A high-rank RNN language model.
\newblock En \emph{International Conference on Learning Representations (ICLR)}, 2018.
\newblock \url{https://arxiv.org/abs/1711.03953}

\bibitem{cs231n}
F.-F. Li \emph{et al.}
\newblock CS231n: Deep Learning for Computer Vision. Notas del curso.
\newblock Stanford University, 2026. «Backpropagation, Intuitions» y «Neural Networks», partes~1 y~2.
\newblock \url{https://cs231n.github.io/}

\bibitem{d2l}
A.~Zhang, Z.~C. Lipton, M.~Li y A.~J. Smola.
\newblock \emph{Dive into Deep Learning}, versión 1.0.3.
\newblock Cambridge University Press, 2023.
\newblock Cap.~5, «Multilayer Perceptrons», y §9.3, «Language Models».
\newblock \url{https://d2l.ai/}

\end{thebibliography}


\end{document}
